% ID: SEL-D0047-OA-E04
% ESTADO: cotejada. Segunda pasada visual y borrador compilado revisados.
% Fuente: Selectividad/sel_2014_matematicas_aplicadas/Matematicas_CCSS_SEPT.pdf
% SHA-256: ee5bd8e23465d55ed58339e7d9adf3d58423ac2966b53d92e39a2f9b35558ab1
% Página: 1; opción A; ejercicio original 4.

\textbf{(2.5 puntos)} La concejalía de Educación de una determinada localidad afirma que el tiempo medio dedicado a la lectura por los jóvenes de entre 15 y 20 años de edad es, a lo sumo, de 8 horas semanales. Para contrastar esta hipótesis, $(H_0:\mu\leq8)$, se escoge al azar una muestra de 100 jóvenes, de entre 15 y 20 años, y se obtiene una media de 8.3 horas de dedicación a la lectura. Supuesto que el tiempo dedicado a la lectura sigue una ley Normal con desviación típica igual a 1 hora, \textquestiondown{}qué se puede decir, a un nivel de significación del 5\%, sobre la afirmación de la concejalía?
